\documentclass[11pt]{article}

% Change "review" to "final" to generate the final (sometimes called camera-ready) version.
% Change to "preprint" to generate a non-anonymous version with page numbers.
\usepackage[review]{acl}

% Standard package includes
\usepackage{times}
\usepackage{latexsym}

% For proper rendering and hyphenation of words containing Latin characters (including in bib files)
\usepackage[T1]{fontenc}
% For Vietnamese characters
% \usepackage[T5]{fontenc}
% See https://www.latex-project.org/help/documentation/encguide.pdf for other character sets

% This assumes your files are encoded as UTF8
\usepackage[utf8]{inputenc}
\usepackage{textgreek}
\usepackage{textcomp}
\usepackage{newunicodechar}
\newunicodechar{₄}{$_4$}
\newunicodechar{【}{[}
\newunicodechar{】}{]}
\usepackage{amsmath}
\usepackage{amssymb}
\usepackage{dsfont}
% This is not strictly necessary, and may be commented out,
% but it will improve the layout of the manuscript,
% and will typically save some space.
\usepackage{microtype}
\usepackage{multirow}

% This is also not strictly necessary, and may be commented out.
% However, it will improve the aesthetics of text in
% the typewriter font.
\usepackage{inconsolata}

%Including images in your LaTeX document requires adding
%additional package(s)
\usepackage{graphicx}
\usepackage[utf8]{inputenc}
\usepackage{amssymb}
% If the title and author information does not fit in the area allocated, uncomment the following
%
%\setlength\titlebox{<dim>}
%
% and set <dim> to something 5cm or larger.

%\title{Agri-R1: Enhancing Agricultural Disease Diagnosis through \\
%Reinforcement Learning with Automated Chain-of-Thought Generation}


\title{Agri-R1: Empowering Generalizable Agricultural Reasoning in \\
Vision-Language Models with Reinforcement Learning}

% Author information can be set in various styles:
% For several authors from the same institution:
% \author{Author 1 \and ... \and Author n \\
%         Address line \\ ... \\ Address line}
% if the names do not fit well on one line use
%         Author 1 \\ {\bf Author 2} \\ ... \\ {\bf Author n} \\
% For authors from different institutions:
% \author{Author 1 \\ Address line \\  ... \\ Address line
%         \And  ... \And
%         Author n \\ Address line \\ ... \\ Address line}
% To start a separate ``row'' of authors use \AND, as in
% \author{Author 1 \\ Address line \\  ... \\ Address line
%         \AND
%         Author 2 \\ Address line \\ ... \\ Address line \And
%         Author 3 \\ Address line \\ ... \\ Address line}

\author{First Author \\
  Affiliation / Address line 1 \\
  Affiliation / Address line 2 \\
  Affiliation / Address line 3 \\
  \texttt{email@domain} \\\And
  Second Author \\
  Affiliation / Address line 1 \\
  Affiliation / Address line 2 \\
  Affiliation / Address line 3 \\
  \texttt{email@domain} \\}

%\author{
%  \textbf{First Author\textsuperscript{1}},
%  \textbf{Second Author\textsuperscript{1,2}},
%  \textbf{Third T. Author\textsuperscript{1}},
%  \textbf{Fourth Author\textsuperscript{1}},
%\\
%  \textbf{Fifth Author\textsuperscript{1,2}},
%  \textbf{Sixth Author\textsuperscript{1}},
%  \textbf{Seventh Author\textsuperscript{1}},
%  \textbf{Eighth Author \textsuperscript{1,2,3,4}},
%\\
%  \textbf{Ninth Author\textsuperscript{1}},
%  \textbf{Tenth Author\textsuperscript{1}},
%  \textbf{Eleventh E. Author\textsuperscript{1,2,3,4,5}},
%  \textbf{Twelfth Author\textsuperscript{1}},
%\\
%  \textbf{Thirteenth Author\textsuperscript{3}},
%  \textbf{Fourteenth F. Author\textsuperscript{2,4}},
%  \textbf{Fifteenth Author\textsuperscript{1}},
%  \textbf{Sixteenth Author\textsuperscript{1}},
%\\
%  \textbf{Seventeenth S. Author\textsuperscript{4,5}},
%  \textbf{Eighteenth Author\textsuperscript{3,4}},
%  \textbf{Nineteenth N. Author\textsuperscript{2,5}},
%  \textbf{Twentieth Author\textsuperscript{1}}
%\\
%\\
%  \textsuperscript{1}Affiliation 1,
%  \textsuperscript{2}Affiliation 2,
%  \textsuperscript{3}Affiliation 3,
%  \textsuperscript{4}Affiliation 4,
%  \textsuperscript{5}Affiliation 5
%\\
%  \small{
%    \textbf{Correspondence:} \href{mailto:email@domain}{email@domain}
%  }
%}

\begin{document}
\maketitle
\begin{abstract}
% Agricultural disease diagnosis via VLMs faces critical barriers: supervised fine-tuning demands massive datasets, produces black-box predictions, and generalizes poorly. While CoT reasoning enhances interpretability and RL improves robustness, existing methods require costly expert annotation or target closed-set medical tasks rather than open-ended agricultural diagnosis.

% We introduce \textbf{Agri-R1}, combining automated CoT generation with GRPO for agricultural VQA. Through vision-language generation and LLM-based filtering, we construct reasoning data from only \textbf{19\%} of samples. Our domain vocabularies with fuzzy-matching reward function handle linguistic diversity in open-ended responses, enabling robust policy learning.

% Our \textbf{3B model achieves competitive performance with 7B~13B baselines}, achieving +23.2\% relative improvement in disease recognition, +33.3\% in knowledge QA, and +26.10-point gain in cross-domain generalization over SFT. Ablations confirm GRPO exploration and CoT reasoning synergistically enhance performance, scaling with question complexity. Agri-R1 \footnote{\url{https://github.com/CPJ-Agricultural/Agri-R1}} establishes a scalable paradigm for data-efficient, interpretable agricultural AI.

%Agricultural disease diagnosis remains a challenge for vision-language models (VLMs), as conventional supervised fine-tuning demands extensive labeled data, yields uninterpretable predictions, and generalizes weakly to real-world scenarios. While reasoning capabilities can enhance model transparency and robustness, existing methods typically rely on costly expert-annotated reasoning traces and are seldom designed for the open-ended, linguistically diverse nature of agricultural queries.

%%% abstract:
Agricultural disease diagnosis challenges VLMs, as conventional fine-tuning requires extensive labels, lacks interpretability, and generalizes poorly. While reasoning improves model robustness, existing methods rely on costly expert annotations and rarely address the open-ended, diverse nature of agricultural queries.
To address these limitations, we propose \textbf{Agri-R1}\footnote{\url{https://github.com/CPJ-Agricultural/Agri-R1}}, a reasoning-enhanced large model for agriculture.
Our framework automates high-quality reasoning data generation via vision-language synthesis and LLM-based filtering, using only 19\% of available samples.
Training employs Group Relative Policy Optimization (GRPO) with a novel proposed reward function that integrates domain-specific lexicons and fuzzy matching to assess both correctness and linguistic flexibility in open-ended responses.
Evaluated on CDDMBench, our resulting 3B-parameter model achieves performance competitive with 7B- to 13B-parameter baselines, showing a +23.2\% relative gain in disease recognition accuracy, +33.3\% in agricultural knowledge QA, and a +26.10-point improvement in cross-domain generalization over standard fine-tuning. Ablation studies confirm that the synergy between structured reasoning data and GRPO-driven exploration underpins these gains, with benefits scaling as question complexity increases.

%The resulting \textbf{3B-parameter model} matches the performance of 7B to 13B baselines, achieving +23.2\% relative improvement in disease recognition, +33.3\% in agricultural knowledge QA, and +26.10 points in cross-domain generalization over standard fine-tuning. Ablations confirm that structured reasoning data and GRPO exploration synergistically drive gains that increase with question complexity.

%Agri-R1\footnote{\url{https://github.com/CPJ-Agricultural/Agri-R1}} offers a scalable, data-efficient paradigm for interpretable and robust reasoning models in agricultural AI.

\end{abstract}

%\noindent\textbf{Code and data:}



\section{Introduction}
\label{sec:introduction}

%Agricultural crop diseases pose a persistent threat to global food security, causing substantial yield losses and economic damage~\cite{savary2020modeling,shahbazi2025losses,fones2020threats,gai2024plant}. Accurate and timely disease diagnosis is critical for effective crop protection, yet it remains challenging due to visual symptom complexity and expertise scarcity in many regions~\cite{upadhyay2025deep,ngugi2024revolutionizing,buja2021advances,mohanty2016using}. Recent advances in Vision-Language Models (VLMs) have shown promise for automated diagnosis through visual question answering (VQA), enabling farmers to submit crop images with natural language queries for diagnostic guidance~\cite{lu2024application,sapkota2025multi}. % ,fones2020threats
Agricultural crop diseases pose a persistent threat to global food security, causing substantial yield losses and economic damage~\cite{savary2020modeling,gai2024plant,shahbazi2025losses}. Accurate and timely diagnosis is essential for effective crop protection, yet remains challenging due to complex visual symptoms and limited expert availability in many regions~\cite{upadhyay2025deep,ngugi2024revolutionizing,buja2021advances,mohanty2016using}. Recent advances in Vision-Language Models (VLMs) show promise for automated diagnosis via visual question answering (VQA), allowing farmers to submit crop images with natural language queries for diagnostic guidance~\cite{lu2024application,sapkota2025multi}.


%The predominant approach to adapting VLMs for agricultural tasks relies on supervised fine-tuning (SFT). While SFT achieves strong in-domain performance, it suffers from three critical limitations that hinder real-world deployment. \textbf{First, data hunger}: SFT requires massive labeled datasets to achieve competitive accuracy, making it costly and impractical for resource-constrained domains~\cite{liu2024multimodal}. \textbf{Second, limited interpretability}: models produce diagnostic labels without explaining their reasoning, creating a ``black-box'' that undermines farmer trust and prevents validation by agricultural extension agents~\cite{zhi2025medgr,chu2025sft}. \textbf{Third, poor domain generalization}: models memorize dataset-specific patterns rather than learning robust diagnostic reasoning, leading to performance degradation when confronted with domain shifts such as new crops, imaging conditions, or co-occurring infections~\cite{pan2025medvlm,wu2023laboratory,nanavaty2024integrating,chen2025sft,li2024getting}.

The dominant paradigm for adapting VLMs to agricultural tasks is supervised fine-tuning (SFT). While effective in-domain, SFT faces three critical limitations that impede real-world deployment: (1) data hunger, requiring massive labeled datasets that are costly to acquire in resource-constrained domains~\cite{liu2024multimodal}; (2) limited interpretability, models produce diagnostic labels without explaining their reasoning, creating a ``black-box'' that undermines farmer trust and prevents validation by agricultural extension agents~\cite{zhi2025medgr,chu2025sft}; and (3) poor generalization, as models memorize dataset-specific patterns rather than robust diagnostic reasoning, leading to sharp performance drops under domain shifts (e.g., new crops, lighting conditions, or co-infections)~\cite{pan2025medvlm,wu2023laboratory,nanavaty2024integrating,chen2025sft}. %,li2024getting
These limitations collectively point to a fundamental gap: the need for models that are not only accurate but also data-efficient, interpretable, and robust to the open-ended diversity of real-world agricultural queries.


%\textbf{Chain-of-Thought (CoT) reasoning}~\cite{lyu2023faithful} offers a path toward interpretability by generating explicit intermediate reasoning steps, while \textbf{reinforcement learning (RL)} presents a compelling alternative to SFT by encouraging models to explore diverse reasoning strategies through reward signals~\cite{shakya2023reinforcement}. Group Relative Policy Optimization (GRPO)~\cite{shao2024deepseekmath,wang2025grpo,tong2025delving} has demonstrated strong generalization in mathematical and coding domains through group-based advantage estimation. However, constructing high-quality CoT data for agriculture is prohibitively expensive, requiring domain experts to manually annotate reasoning chains. Moreover, existing RL applications in medical VQA~\cite{yi2022automated,hu2023reinforcement} primarily target \emph{closed-set multiple-choice questions} with binary rewards, whereas agricultural VQA demands \emph{open-ended responses} that exhibit high linguistic diversity—a challenge unaddressed by prior work.
Structured reasoning enhances model transparency by generating explicit intermediate steps, while reinforcement learning (RL) offers an alternative to SFT by promoting diverse reasoning strategies through reward guidance~\cite{shakya2023reinforcement}. Group Relative Policy Optimization (GRPO)~\cite{shao2024deepseekmath,wang2025grpo,tong2025delving} has achieved strong generalization in mathematical and coding tasks via group-based advantage estimation. However, a direct application in agriculture faces two synergistic bottlenecks. First, constructing high-quality CoT data is prohibitively expensive, requiring domain experts to manually annotate reasoning chains. Second, existing RL applications in medical VQA~\cite{yi2022automated,hu2023reinforcement} primarily target \emph{closed-set multiple-choice questions} with binary rewards. 
This paradigm is fundamentally mismatched with agricultural VQA, which requires evaluating open-ended, linguistically diverse responses for both factual correctness and reasoning quality—a challenge that remains unaddressed in prior work.
%This paradigm is fundamentally mismatched with agricultural VQA, which demands \emph{open-ended responses} that exhibit high linguistic diversity and require nuanced evaluation of both factual correctness and reasoning quality—a critical challenge unaddressed by prior work.

%However, creating high-quality reasoning data for agriculture demands expensive expert annotation of step-by-step traces. Moreover, existing RL methods in medical VQA~\cite{yi2022automated,hu2023reinforcement} focus on closed-set multiple-choice questions with binary rewards, failing to address the open-ended, linguistically diverse responses required in agricultural VQA—where farmers' natural queries exhibit high variability.


%In this work, we introduce \textbf{Agri-R1}, the first GRPO-based framework for open-ended agricultural VQA. As illustrated in Figure~\ref{fig:framework}, our approach integrates three key innovations to address the aforementioned challenges. \textbf{First}, we eliminate manual CoT annotation costs through an automated pipeline that generates reasoning chains via vision-language models and filters high-quality data using LLM-as-a-Judge, constructing a compact dataset from only \textbf{19\%} of the original corpus. \textbf{Second}, we construct domain vocabularies and design a fuzzy-matching reward function to handle the linguistic diversity inherent in open-ended agricultural responses—fundamentally distinct from binary-reward closed-set medical VQA systems. \textbf{Third}, we demonstrate that GRPO-driven policy optimization enables our compact \textbf{3B-parameter model} to achieve superior accuracy and cross-domain generalization compared to larger baselines trained on full datasets.

To overcome these bottlenecks, we introduce \textbf{Agri-R1}, the first GRPO-based framework designed specifically for open-ended, reasoning-enhanced agricultural VQA. We integrates three key innovations to simultaneously achieve data efficiency, interpretability, and robustness
%. As illustrated in Figure~\ref{fig:framework}
: (1) we eliminate manual CoT annotation costs through an automated pipeline that synthesizes reasoning chains via VLMs and filters high-quality data using LLM-as-a-Judge, constructing a compact yet powerful dataset from only 19\% of the original corpus; (2) to address the unique challenge of evaluating open-ended answers, we construct agricultural domain vocabularies and design a novel fuzzy-matching reward function. This function assesses not just correctness but also the linguistic appropriateness of responses, enabling effective policy optimization far beyond binary rewards; (3) we demonstrate that GRPO-driven policy optimization, fueled by our automated reasoning data and specialized reward, enables a remarkably compact \textbf{3B-parameter model} to achieve superior accuracy and cross-domain generalization compared to significantly larger baselines trained on full datasets.

Our primary contributions are as follows:
\begin{itemize}
    \item We propose Agri-R1, the first GRPO-based framework designed for agricultural disease diagnosis, introducing an automated pipeline to generate and filter reasoning data without relying on expert annotations.
    \item We design a novel reward mechanism based on agricultural lexicons and fuzzy matching to evaluate both correctness and flexibility in open-ended agricultural responses, addressing a critical limitation of binary-reward systems.
    \item We show that a significantly smaller model trained with our framework surpasses larger baselines in accuracy, reasoning ability, and cross-domain generalization, demonstrating the synergistic effect of structured reasoning data and reinforcement learning exploration.
\end{itemize}


% Our contributions are as follows:
% \begin{itemize}
%     \item \textbf{First GRPO framework for open-ended agricultural VQA}, with an automated CoT construction pipeline combining vision-language generation and quality filtering, eliminating costly manual annotation.
%     \item \textbf{Domain vocabularies with fuzzy-matching reward function} addressing linguistic diversity in open-ended responses—a key challenge distinct from binary-reward closed-set medical VQA, enabling robust policy optimization.
%     \item \textbf{Strong empirical results demonstrating dual efficiency}: our 3B model outperforms 7B baselines while training on only 19\% data, achieving +23.2\% relative improvement in disease recognition, +33.3\% in knowledge QA, and +26.10-point cross-domain generalization gain over SFT; ablation studies reveal synergistic contributions of GRPO exploration and CoT reasoning, with CoT's contribution scaling with question complexity.
% \end{itemize}

%Agri-R1 establishes a new paradigm for data-efficient, interpretable agricultural AI through automated reasoning data construction and reward-guided policy optimization.


%%%%%%%%%%%%%%%%%%%%%%%%%%%%%%%%%%%%%%%%%%%%%%%%%%%%%%%%%%
\begin{figure*}[t]
\centering
\includegraphics[width=\textwidth]{figure1_construct image.jpg}
\caption{\textbf{Overview of our Two-Stage GRPO Framework for Agricultural Disease Reasoning.} Stage 1 transforms raw VQA pairs into reasoning exemplars: DeepSeek-VL2 generates reasoning chains, GPT-4 filters outputs (threshold $\tau$=8.0/10.0). Stage 2 employs GRPO-based policy learning with Domain Vocabulary Construction: 5-tier fuzzy matching handles linguistic diversity, three-component reward function (Format + Answer + Reasoning) guides optimization, and group relative advantage normalization (n=3 samples) enables stable updates. This pipeline enables our 3B model to learn robust reasoning from synthesized data.}
\label{fig:framework}
\end{figure*}
%%%%%%%%%%%%%%%%%%%%%%%%%%%%%%%%%%%%%%%%%%%%%%%%%%%%%%%

\section{Related Work}
\label{sec:related_work}

\paragraph{Agricultural Vision-Language Models}
Recent advances in vision-language models (VLMs) have spurred domain-specific adaptations for agricultural disease diagnosis~\cite{zhou2024few,awais2025agrogpt,arshad2025leveraging}. Existing studies typically follow two paradigms: some focus on compact model design, such as ~\citet{cao2025small}, who employ image-text contrastive learning for few-shot crop disease identification; others incorporate domain knowledge for enhanced representational alignment, such as ~\citet{yao2024multimodal}, who integrate meteorological indicators for multimodal drought detection. Large-scale data initiatives, such as AGBase-2000K, have further facilitated knowledge integration through comprehensive multimodal agricultural corpora~\cite{gauba2025agmmu}. Despite these efforts, ~\citet{liu2024multimodal} reveal that models trained via supervised fine-tuning (SFT) remain prone to performance degradation under domain shifts, reflecting the limited robustness and interpretability of current approaches in open-ended agricultural VQA scenarios.


%Recent advances in vision-language models (VLMs) have spurred domain-specific adaptations for agricultural disease diagnosis~\cite{zhou2024few,awais2025agrogpt,arshad2025leveraging}. ~\citet{cao2025small} employ image-text contrastive learning for small-sample cucumber disease identification, achieving high accuracy. ~\cite{yao2024multimodal} propose S-DNet, which fuses images with meteorological indices for winter wheat drought detection, reaching 96.4\% accuracy. Large-scale efforts include AGBase-2000K, a comprehensive agricultural multimodal knowledge base that improves model performance by 4.7\%~\cite{gauba2025agmmu}. Liu et al. construct a 35,000-sample expert-guided benchmark, revealing that SFT-trained models struggle with out-of-distribution generalization—a critical limitation for real-world agricultural deployment.

% While these methods achieve strong in-domain performance, they predominantly rely on supervised fine-tuning (SFT), which tends to overfit narrow training distributions and provides limited transparency in diagnostic reasoning. Our work addresses these limitations by adopting reinforcement learning (RL) to encourage generalizable reasoning without requiring extensive domain-specific annotations.

\paragraph{Chain-of-Thought for Interpretability}
The inherent ``black-box'' nature of large language models presents a fundamental barrier to their adoption in high-stakes applications such as agriculture, where the need for transparent and trustworthy decision-making is paramount~\cite{sun2022black,bommasani2021opportunities,martin2024xai}. To address this, CoT prompting~\cite{wei2022chain} has emerged as a prominent method for enhancing model interpretability, eliciting explicit, step-by-step reasoning paths from models. Subsequent work has aimed to improve the reliability of CoT; for instance, ~\citet{wang2022self} enhance CoT robustness through self-consistency by aggregating predictions across multiple reasoning paths. However, a critical bottleneck persists: the manual curation of high-quality, domain-specific CoT demonstrations remains prohibitively expensive and difficult to scale~\cite{wang2025beyond,lightman2023letsverifystepstep,kim2023cotcollectionimprovingzeroshot}. This challenge is especially pronounced in agriculture, where expert knowledge is required to validate the correctness and relevance of diagnostic reasoning chains, underscoring the need for scalable, automated solutions for CoT generation.

% The inherent ``black-box'' nature of large language models poses significant challenges for high-stakes applications, particularly in agriculture where practitioners require transparent decision-making processes~\cite{sun2022black,bommasani2021opportunities,martin2024xai}. Chain-of-thought (CoT) prompting~\cite{wei2022chain} has emerged as a prominent technique for eliciting interpretable reasoning by embedding intermediate reasoning steps in few-shot examples. Wang et al. further enhance CoT robustness through self-consistency, aggregating predictions across multiple reasoning paths~\cite{wang2022self}. However, manually curating high-quality CoT demonstrations for specialized domains is prohibitively expensive and difficult to scale~\cite{wang2025beyond,lightman2023letsverifystepstep,kim2023cotcollectionimprovingzeroshot}.
% Unlike prior work that relies on costly expert annotations or teacher model distillation, our approach leverages RL to automatically discover interpretable reasoning patterns during training.

\paragraph{Reinforcement Learning for VLM Post-Training}
Reinforcement learning (RL) provides a powerful paradigm for aligning models with desired behaviors through reward feedback, offering an alternative to supervised fine-tuning that emphasizes outcome-driven optimization~\cite{christiano2017deep,ladosz2022exploration}. Within this paradigm, Proximal Policy Optimization (PPO)~\cite{schulman2017proximal} has established itself as a stable and effective algorithm for policy learning in language and vision-language settings. Recent advancements have sought to improve RL efficiency and scalability. Group Relative Policy Optimization (GRPO)~\cite{shao2024deepseekmath} simplifies the training architecture by replacing the learned value function with group-based advantage estimation, significantly reducing computational cost while maintaining stable convergence. This approach has demonstrated strong reasoning capabilities in domains such as mathematics~\cite{shao2024deepseekmath} and coding~\cite{guo2025deepseek}. Similarly, in specialized fields like medical vision-language understanding, RL has been successfully adapted to address data scarcity and improve cross-modal generalization, as seen in works like ~\citet{zhi2025medgr}.

%However, reinforcement learning—particularly GRPO—has yet to be applied to open-ended agricultural vision-language reasoning, marking a significant unexplored frontier. Agricultural VQA uniquely requires interpretable step-by-step diagnostic reasoning amid high linguistic variability, severe data scarcity, and pronounced domain shifts—challenges ill-suited to existing RL methods reliant on binary or closed-set rewards.
%We pioneer the first GRPO framework for this domain, demonstrating that reward-guided optimization yields breakthrough gains in accuracy, generalization, transparency, and robustness, all without costly expert-annotated reasoning traces.

The application of RL to open-ended agricultural VQA remains an open challenge, with no prior work adapting GRPO to this domain. Unlike medical or mathematical tasks, agricultural diagnosis demands interpretable reasoning under linguistic diversity, data scarcity, and domain shifts. Existing RL methods rely on binary or multiple-choice rewards, which fail to evaluate free-form, agriculturally grounded explanations. To our knowledge, we present the first GRPO-based framework for agricultural VQA, introducing a domain-aware reward design and automated reasoning data synthesis to jointly improve accuracy, generalization, and interpretability—without expert annotations.


%Reinforcement learning offers a compelling alternative to SFT for aligning models with desired behaviors through reward signals rather than direct supervision~\cite{christiano2017deep,ladosz2022exploration}. Proximal Policy Optimization (PPO)~\cite{schulman2017proximal} has become a foundational algorithm for policy learning in language models, balancing exploration with stable convergence. More recently, Group Relative Policy Optimization (GRPO)~\cite{shao2024deepseekmath} eliminates the need for learned value functions by estimating advantages via group-based reward normalization, significantly reducing computational overhead while maintaining training stability. DeepSeek-R1~\cite{guo2025deepseek} demonstrates that GRPO can incentivize sophisticated reasoning in mathematical and coding domains. In the medical domain, MedGR²~\cite{zhi2025medgr} applies RL with dynamic reward modeling to overcome data scarcity, achieving robust cross-modality generalization.




% However, these RL-based approaches remain largely unexplored in agricultural VQA, where data scarcity, domain shift, and the need for interpretable reasoning converge.
% Our work bridges this gap by adapting GRPO to agricultural disease diagnosis, demonstrating that RL-driven post-training can simultaneously enhance accuracy, generalization, and transparency without requiring expert-annotated reasoning chains.


%Reinforcement learning offers a compelling alternative to SFT for aligning models with desired behaviors through reward signals rather than direct supervision~\cite{christiano2017deep,ladosz2022exploration}. Proximal Policy Optimization (PPO)~\cite{schulman2017proximal} has become a foundational algorithm for policy learning in language models, balancing exploration with stable convergence. More recently, Group Relative Policy Optimization (GRPO)~\cite{shao2024deepseekmath} eliminates the need for learned value functions by estimating advantages via group-based reward normalization, significantly reducing computational overhead while maintaining training stability. DeepSeek-R1~\cite{guo2025deepseek} demonstrates that GRPO can incentivize sophisticated reasoning in mathematical and coding domains. In the medical domain, MedGR²~\cite{zhi2025medgr} applies RL with dynamic reward modeling to overcome data scarcity, achieving robust cross-modality generalization.


\section{Methodology}
\label{sec:methodology}
%To address data scarcity and reasoning generalization in agricultural VQA, we propose a two-stage framework combining generative CoT enhancement with GRPO-based reinforcement learning (Figure~\ref{fig:framework}). Stage 1 constructs high-quality reasoning data through automated generation and GPT-4 filtering. Stage 2 employs GRPO with domain-specific reward functions to train a robust policy. This paradigm shifts the bottleneck from manual annotation to intelligent data synthesis.

Figure~\ref{fig:framework} illustrates our framework. It consists of two stages: a Generative Reasoning Enhancement Engine for constructing a high-quality reasoning dataset, followed by a GRPO Reinforcement Learning stage for training a robust policy with domain-specific rewards.



%\subsection{Chain-of-Thought Data Augmentation}
\subsection{Generative Reasoning Enhancement}
\label{sec:cot_augmentation}

To enable interpretable reasoning without manual annotation, we employ the three-step pipeline in Figure~\ref{fig:framework} (STAGE1): (1) \textbf{Data Processing}—resize images; (2) \textbf{Reasoning Data Generation}—DeepSeek-VL2~\cite{wu2024deepseek} generates reasoning chains in structured format $\langle\text{think}\rangle R \langle/\text{think}\rangle \langle\text{answer}\rangle A \langle/\text{answer}\rangle$; (3) \textbf{Judge and Enhanced}—GPT-4 filters reasoning quality (threshold $\tau$=8.0/10.0), with low-scoring chains regenerated via feedback-guided prompting. (Details are provided in Appendix~\ref{sec:appendix_reasoning}).

% Detailed synthesis pipeline and statistics are provided in Appendix~\ref{sec:appendix_reasoning}.

\subsection{GRPO Reinforcement Learning}
\label{sec:grpo}

\subsubsection{Group Relative Policy Optimization}
%Unlike approaches requiring neural reward models, GRPO directly optimizes the policy using group-based advantage estimation~\cite{zheng2025group}. Given a policy $\pi_\theta$ (initialized from supervised fine-tuning) and a reference policy $\pi_{\text{ref}}$, GRPO samples $G$ candidate responses for each input pair $(I, q)$:

GRPO optimizes the policy $\pi_\theta$ using group-based advantage estimation~\cite{zheng2025group}, without requiring a separate reward model. For each input $(I, q)$, we sample $G$ responses:

\begin{equation}
o_i \sim \pi_\theta(\cdot \mid I, q), \quad i = 1, \ldots, G
\label{eq:sample}
\end{equation}

\noindent where $o_i$ is a candidate response, $I$ the input image, and $q$ the question. Each response receives a scalar reward $r_i$ from our reward function (Section~\ref{sec:reward}).
%where $o_i$ represents the $i$-th candidate response generated by the current policy, $I$ is the input image, and $q$ is the question. Each response receives a scalar reward $r_i$ from our reward function (Section~\ref{sec:reward}).

\noindent The group relative advantage normalizes rewards within each group of responses to stabilize learning:
%The group relative advantage normalizes rewards within each sample group to address inter-sample reward variance:

\begin{equation}
\begin{aligned}
A_i &= \frac{r_i - \mu_G}{\sigma_G + \epsilon}, \quad
\mu_G = \frac{1}{G}\sum_{j=1}^G r_j, \\
\sigma_G &= \sqrt{\frac{1}{G}\sum_{j=1}^G (r_j - \mu_G)^2}
\end{aligned}
\label{eq:advantage}
\end{equation}

\noindent where $A_i$ is the advantage for candidate $i$, $\mu_G$ and $\sigma_G$ are the group’s mean and standard deviation, and $\epsilon$ is a small constant for stability. This normalization helps the model learn from relative quality differences within each group.

The GRPO objective balances policy improvement with KL regularization:
% \noindent where $A_i$ is the advantage for candidate $i$, $\mu_G$ is the mean reward, $\sigma_G$ is the standard deviation, and $\epsilon$ is a small constant for numerical stability. This normalization focuses learning on within-group quality differences rather than absolute reward magnitudes.

% The GRPO objective function combines a clipped policy gradient term with KL divergence regularization:

\begin{equation}
\begin{aligned}
\mathcal{J}_G(\theta) &= \mathbb{}b{E}_{(I,q)\sim\mathcal{D}} \Bigg[
\frac{1}{G} \sum_{i=1}^G \min\big( \rho_i A_i, \\
&\quad \text{clip}(\rho_i, 1-\varepsilon, 1+\varepsilon) A_i \big) \\
&\quad - \beta \cdot D_{KL}(\pi_\theta \| \pi_{ref}) \Bigg]
\end{aligned}
\label{eq:grpo_objective}
\end{equation}

\noindent where \(\mathcal{J}_G\) is the GRPO objective; \(\rho_i\) is the probability ratio between current and old policies; clip and a $KL$ penalty enforce conservative policy updates.
%where $\mathcal{J}_G$ is the GRPO objective; $\rho_i = \pi_\theta(o_i \mid I, q) / \pi_{\theta_{old}}(o_i \mid I, q)$ is the probability ratio; $\text{clip}(\rho_i, 1-\varepsilon, 1+\varepsilon)$ restricts updates; $\varepsilon$ is the clipping threshold; $\beta$ controls KL penalty strength; and $D_{KL}(\pi_\theta \| \pi_{ref})$ measures policy divergence.

\subsubsection{Reward Function Design}
\label{sec:reward}

A key challenge in agricultural VQA is designing reward functions for \textbf{open-ended} responses with high linguistic diversity~\cite{qian2025toolrl,eschmann2021reward,liu2024multimodal,lai2025med,pan2025medvlm}. We construct domain-specific vocabularies $\mathcal{V}_p$ and $\mathcal{V}_d$ for synonym recognition, then define a three-component reward function:

\begin{equation}
R(o) = w_f R_f(o) + w_a R_a(o) + w_r R_r(o)
\label{eq:total_reward}
\end{equation}

\noindent where $o$ is the candidate response; $R_f$, $R_a$, and $R_r$ denote Format, Answer Exact Match, and Reasoning Quality rewards respectively; $w_f=0.5$ (17\%), $w_a=2.0$ (67\%), and $w_r=0.5$ (17\%) are the component weights; and $R(o) \in [0, 3.0]$ is the total reward. Detailed scoring criteria are provided in Appendix~\ref{sec:appendix_reward}.

\paragraph{Domain Vocabularies.}
We construct domain-specific vocabularies $\mathcal{V}_p$ (plant species) and $\mathcal{V}_d$ (disease types) from CDDMBench's 15 crop types and 20 disease categories. Each entry includes canonical names, scientific nomenclature (e.g., ``tomato'' $\leftrightarrow$ ``\emph{Solanum lycopersicum}''), and colloquial variations to handle linguistic diversity in agricultural diagnosis (Appendix~\ref{sec:appendix_reward} provides the detailed vocabulary construction).



\paragraph{Format Reward.}
This component ensures structured output with required tags and quality metrics:

\begin{equation}
R_f(o) = \begin{cases}
\sum_{c \in C_f} w_c \cdot r_c(o) & \text{if tags exist} \\
0 & \text{otherwise}
\end{cases}
\label{eq:format_reward}
\end{equation}

\noindent where $C_f = \{\text{struct}, \text{steps}, \text{content}, \text{length}, \text{quality}\}$ evaluates basic structure with $\langle\text{think}\rangle$ (reasoning) $\langle/\text{think}\rangle$ and $\langle\text{answer}\rangle$ (response) $\langle/\text{answer}\rangle$ tags ($w=0.15$), step structure and content quality ($w=0.15, 0.10$), and appropriate think/answer lengths ($w=0.05$ each), summing to 0.5.

\paragraph{Answer Keyword Reward.}
This component evaluates diagnostic accuracy using the domain vocabularies. For diagnostic questions, we employ weighted dual matching:

\begin{equation}
\begin{aligned}
R_a^{\text{diag}}(o) &= w_p \cdot M_p(o, a) \\
&\quad + w_d \cdot M_d(o, a)
\end{aligned}
\label{eq:answer_diag}
\end{equation}

\noindent where $w_p=0.8$ and $w_d=1.2$ weight plant and disease matching; $M_p(o, a)$ and $M_d(o, a)$ measure matches using five-tier fuzzy scoring (1.0 exact, 0.85 high-quality, 0.7 partial, 0.5 keyword, 0.25 weak relevance).

For prevention/control questions, we match against method categories:

\begin{equation}
R_a^{\text{ctrl}}(o) = \sum_{c} w_c \cdot \mathds{1}[\text{Contains}(o, \mathcal{V}_c)]
\label{eq:control_reward}
\end{equation}

\noindent where $c \in \{ch, cu, b, t\}$ denotes chemical ($w=0.6$), cultural ($w=0.5$), biological ($w=0.5$), and timing ($w=0.4$) methods; $\mathcal{V}_c$ are category vocabularies; $\mathds{1}[\cdot]$ indicates keyword presence.

\paragraph{Reasoning Quality Reward.}
This component evaluates Chain-of-Thought quality through three dimensions:

\begin{equation}
R_r(o) = \sum_{d \in D_r} w_d \cdot r_d(o)
\label{eq:reasoning_reward}
\end{equation}

\noindent where $D_r = \{\text{logic}, \text{prof}, \text{comp}\}$ evaluates logical coherence through causal patterns (e.g., "observe...because") and step connections ($w=0.25$), professional terminology usage in appropriate diagnostic context ($w=0.15$), and reasoning chain completeness covering observation→analysis→conclusion flow ($w=0.10$). Detailed scoring criteria are provided in Appendix~\ref{sec:appendix_reward}.

\paragraph{Dynamic Evaluation.}
Our reward function dynamically selects evaluation criteria based on question type, distinguishing between diagnostic queries and control questions via separate scoring formulations in Equation~\ref{eq:answer_diag} and Equation~\ref{eq:control_reward}. To address the inherent lexical variation in open-ended agricultural responses, the function incorporates a five-tier fuzzy matching mechanism that ranges from exact match to weak relevance~\cite{reichard2025open}. The complete reward pipeline processes each candidate response $o_i$ through stages of format validation, keyword extraction, and semantic matching, culminating in a final scalar reward $r_i$ within the bounded interval $[0, 3.0]$ for subsequent GRPO optimization as defined in Equation~\ref{eq:grpo_objective}.

% Our reward function dynamically selects evaluation criteria based on question type (Equation~\ref{eq:answer_diag} for diagnosis vs. Equation~\ref{eq:control_reward} for control) and employs five-tier fuzzy matching (ranging from exact match to weak relevance) to handle lexical variation in open-ended agricultural responses~\cite{reichard2025open}.

% The complete reward pipeline processes each candidate response $o_i$ through format validation, keyword extraction, and fuzzy matching, producing $r_i \in [0, 3.0]$ for GRPO training (Equation~\ref{eq:grpo_objective}).

% ==============================================================================
% Optimized Experiments and Results Sections for ACL Paper
% Strictly following ACL LaTeX template format
% Updated with new tables (Liu et al. 2024, Zhang et al. 2024)
% Target: 4.5-5 pages including figures and tables
% ==============================================================================

% ==============================================================================
% Optimized Experiments and Results Sections for ACL Paper
% Strictly following ACL LaTeX template format
% Updated with new tables (Liu et al. 2024, Zhang et al. 2024)
% Target: 4.5-5 pages including figures and tables
% ==============================================================================



\section{Experiments}
\label{sec:experiments}

\subsection{Datasets and Evaluation}

\paragraph{Dataset Construction.}
We build training datasets based on CDDMBench~\cite{liu2024multimodal}. The SFT training set uses the full CDDMBench dataset (1.05M samples) in standard VQA format. The GRPO training set applies stratified sampling to obtain 200,005 samples (19\%), preserving class distribution across 15 crops and 20 diseases. This ratio aligns with reward-guided data efficiency findings~\cite{zhi2025medgr}. These samples undergo automated reasoning synthesis via DeepSeek-VL2 generation and GPT-4 quality validation ($\tau$=8.0/10.0), producing reasoning-enhanced data with structured \texttt{<think>} and \texttt{<answer>} tags. Detailed synthesis pipeline and statistics are provided in Appendix~\ref{sec:appendix_reasoning}.


\paragraph{Evaluation Protocol.}
Following CDDMBench protocol, we evaluate on: (1) In-distribution test set (3,963 samples) using keyword matching accuracy for crop/ disease recognition; (2) Disease Knowledge QA (20 samples) scored by GPT-4 (0-10 scale) on professionalism, completeness, and practicality; (3) We also evaluate on AgMMU benchmark (770 samples)~\cite{gauba2025agmmu} for cross-scenario generalization using harmonic mean across five subtasks.


%%% table1 main_res_CDDMBench %%%%%%%%%%%%%%%%%%%%%%%%%%%%%%%%%
\input{t1_main_res_CDDMBench}
%%%%%%%%%%%%%%%%%%%%%%%%%%%%%%%%%%%%

\subsection{Training Configuration}

We adopt Qwen2.5-VL-3B-Instruct~\cite{bai2025qwen2} as our base VLM. Training is conducted on 4 NVIDIA A800 80GB GPUs with DeepSpeed ZeRO-3 optimization. The hyperparameters include a batch size of 160, AdamW optimizer with learning rate $8\times10^{-7}$ and cosine schedule warmup, gradient clipping at 0.3, and BF16 mixed precision. The model is trained for 3 epochs, with the optimal checkpoint selected at step 1,800. For GRPO training, we sample $K=3$ candidate responses per query with temperature $T=0.7$. The reward function evaluates format compliance, answer accuracy, and reasoning quality, with the KL divergence stabilizing between 0.036 and 0.040. The total training time is 98 hours. Full implementation details are provided in Appendix~\ref{sec:empirical_config}.

% We employ Qwen2.5-VL-3B-Instruct~\cite{bai2025qwen2} as base model. Training uses 4 NVIDIA A800 80GB GPUs with DeepSpeed ZeRO-3. Key settings: batch size 160 (10/GPU × 4 accum × 4 GPUs), AdamW (lr=8×10$^{-7}$), cosine schedule (15\% warmup), gradient clipping 0.3, BF16 precision, 3 epochs, optimal checkpoint at step 1,800. For GRPO training: Generates K=3 candidates (T=0.7). Reward function encompasses format compliance, answer accuracy (67\%), and reasoning quality (17\%). KL divergence stable at 0.036-0.040. Training time: 98 hours. Detailed configuration in Appendix~\ref{sec:empirical_config}.

\subsection{Baselines}

We evaluate our method against the following baselines: \textbf{Zero-shot:} Using the pretrained Qwen2.5-VL-3B-Instruct model with only task prompts. \textbf{Few-shot:} The zero-shot approach augmented with 5 in-context examples. \textbf{SFT:} Supervised fine-tuning on the complete CDDMBench dataset (1.05M samples). \textbf{GRPO:} A reinforcement learning variant optimized with answer correctness rewards only, without explicit reasoning. \textbf{Reasoning-Enhanced GRPO (Ours):} Our complete two-stage framework, which integrates automated reasoning data synthesis and reasoning-aware reward functions.
Furthermore, we compare our results to published baselines, namely: \textbf{CDDMBench}~\cite{lu2024application}: A method based on Supervised Fine-Tuning with LoRA, applied to crop disease datasets. \textbf{CPJ}~\cite{zhang2025cpjexplainableagriculturalpest}: A training-free approach that utilizes explainable captions and employs an LLM-as-Judge for evaluation.

% Zero-shot: Pretrained Qwen2.5-VL-3B-Instruct with prompts only.
% Few-shot: Zero-shot + 5 in-context examples.
% SFT: Supervised fine-tuning on full CDDMBench (1.05M samples).
% GRPO: RL with answer correctness rewards only (no explicit reasoning).
% Reasoning-Enhanced GRPO (Ours): Full framework with automated reasoning synthesis + reasoning-aware rewards.

% Published Baselines: We compare against CDDMBench~\cite{lu2024application}, which employs SFT with LoRA on crop disease datasets, and CPJ~\cite{zhang2025cpjexplainableagriculturalpest}, a training-free approach using explainable captions and LLM-as-Judge evaluation.



% ==============================================================================
% Section 4: Results
% ==============================================================================




\subsection{Main Results}
\label{sec:results}

\paragraph{Overall Performance on CDDMBench}
Table~\ref{tab:cddmbench_comparison} presents comprehensive results comparing our approach with published baselines. % Despite using only 19\% training data, our Reasoning-Enhanced GRPO model achieves competitive performance across all three evaluation tasks. 
We present three key observations from our evaluation: 
\textbf{1). Crop Recognition:} Reasoning-Enhanced GRPO (92.58\%) achieves +1.61\% absolute gain over SFT (90.97\%), with GRPO contributing +1.36\% and explicit reasoning adding +0.25\%. 
\textbf{2). Disease Recognition:} Reasoning-Enhanced GRPO (72.50\%) achieves +23.2\% relative gain over SFT, with GRPO contributing +10.59\% and reasoning adding +3.07\%. RL-based exploration (GRPO) provides the dominant improvement, while explicit reasoning enhances fine-grained symptom differentiation. 
\textbf{3). Knowledge QA:} Reasoning-Enhanced GRPO (84.0) matches state-of-the-art Gpt-5-Nano approaches with +33.3\% relative gain over SFT. Critically, reasoning's contribution (+11.51 points) exceeds GRPO's contribution (+9.49 points) on this task, confirming that explicit reasoning chains are essential for multi-step knowledge integration—not merely for transparency, but for fundamental capability enhancement.

% Key Observations:
% \begin{itemize}
% \item Crop Recognition: Reasoning-Enhanced GRPO (92.58\%) achieves +1.61\% absolute gain over SFT (90.97\%), with GRPO contributing +1.36\% and explicit reasoning adding +0.25\%.

% \item Disease Recognition: Reasoning-Enhanced GRPO (72.50\%) achieves +23.2\% relative gain over SFT, with GRPO contributing +10.59\% and reasoning adding +3.07\%. RL-based exploration (GRPO) provides the dominant improvement, while explicit reasoning enhances fine-grained symptom differentiation.

% \item Knowledge QA: Reasoning-Enhanced GRPO (84.0) matches state-of-the-art Gpt-5-Nano approaches with +33.3\% relative gain over SFT. Critically, reasoning's contribution (+11.51 points) exceeds GRPO's contribution (+9.49 points) on this task, confirming that explicit reasoning chains are essential for multi-step knowledge integration—not merely for transparency, but for fundamental capability enhancement.
% \end{itemize}

\begin{figure}[!t]
\centering
\includegraphics[width=0.96\columnwidth]{figure3_agmmu_radar_improved.pdf}
\caption{AgMMU task distribution. Reasoning-Enhanced GRPO (red) outperforms SFT (blue) on visual tasks while showing balanced performance.}
\label{fig:agmmu_radar}
\end{figure}

\paragraph{Generalization Performance on AgMMU-MCQs}

We evaluate generalization capability on AgMMU-MCQs, a subset of the AgMMU benchmark testing agricultural reasoning across five tasks. Reasoning-Enhanced GRPO (66.10\%) matches LLaVA-1.5-13B (66.73\%) and surpasses Qwen-VL-7B (62.34\%) and Claude 3 Haiku (62.00\%) with only 3B parameters. Figure~\ref{fig:agmmu_radar} visualizes performance across five tasks. SFT's performance drops from CDDMBench (90.97\%) to AgMMU-MCQs (40.00\%)-—a 50.97-point collapse. In contrast, GRPO (no explicit reasoning) maintains 59.75\% despite identical 19\% training data, demonstrating +19.75-point better generalization. Reasoning-Enhanced GRPO extends this to +26.10 points, confirming that GRPO's exploration learns domain-invariant features. SFT's degradation on visual tasks contrasts with GRPO's robust performance, highlighting RL's transferable representations.



%%%%%%%%%%%%%%%%%%%%%%%%%%%%%%%%%%%%%%%%%%%%%%%%%%%%%%
\begin{table}[!t]
\centering
% \small
\resizebox{.48\textwidth}{!}{
\begin{tabular}{lrcc}
\hline
\textbf{Crop} & \textbf{Freq.} & \textbf{SFT} & \textbf{Reasoning-GRPO} \\
\hline
\multicolumn{4}{l}{\textit{High-freq. (>5\%) -- Stable ($\sigma$=3.2\%)}} \\
%\\
Tomato & 37.19\% & 90.95\% & \textbf{96.05\%} (+5.10) \\
Apple & 29.48\% & 90.94\% & \textbf{97.69\%} (+6.75) \\
Corn & 8.35\% & 91.12\% & \textbf{95.87\%} (+4.75) \\
\hline
\multicolumn{4}{l}{\textit{Mid-freq. (2-5\%) -- Moderate ($\sigma$=8.7\%)}} \\
%\\
Potato & 4.21\% & 90.88\% & \textbf{94.23\%} (+3.35) \\
Pepper & 3.15\% & 91.05\% & \textbf{93.87\%} (+2.82) \\
\hline
\multicolumn{4}{l}{\textit{Low-freq. (<2\%) -- High Variance ($\sigma$=24.5\%)}} \\
%\\
Grape & 1.28\% & 90.84\% & \textbf{100.00\%} (+9.16) \\
Cherry & 1.37\% & 91.30\% & 31.88\% (\textcolor{red}{-59.42}) \\
Strawberry & 1.18\% & 90.72\% & 86.54\% (-4.18) \\
\hline
\end{tabular}}
\caption{Crop recognition by training frequency. Low-freq. crops show high variance ($\sigma$=24.5\%), while high-freq. crops exhibit stable improvements ($\sigma$=3.2\%).}
\label{tab:crop_analysis}
\end{table}
%%%%%%%%%%%%%%%%%%%%%%%%%%%%%%%%%%%%%%%%%%%%%%%%%%%%%%




\section{Analysis}
%\subsection{Detailed Analysis: Crop Recognition Variance}
\subsection{Frequency-Induced Bias in Crop Recognition}
Table~\ref{tab:crop_analysis} reveals frequency-dependent performance variance in crop recognition, with stability quantified by standard deviation ($\sigma$). High-frequency crops (>5\%) show consistent gains ($\sigma$=3.2\%), whereas low-frequency crops (<2\%) exhibit extreme variability ($\sigma$=24.5\%), ranging from substantial improvements (e.g., +9.16\% on Grape) to severe degradation (e.g., -59.42\% on Cherry). This collapse in low-frequency crops stems from gradient competition: dominant classes like Apple (29.48\% frequency) receive disproportionately more updates (~21× compared to Cherry at 1.37\%), overwriting rarer representations. Shared taxonomic features (e.g., Rosaceae family) amplify the issue, as Apple-specific patterns spuriously dominate category embeddings. Ultimately, this exposes a core limitation of unweighted group-relative advantage estimation (Equation~\ref{eq:advantage}), where high-frequency samples bias batch statistics and hinder long-tail robustness.

%%%%%%%%%%%%%%%%%%%%%%%%%%%%%%%%%%%%%%%%%%%%%%%%%%%%%%%
\begin{figure}[!t]
\centering
\includegraphics[width=\columnwidth]{figure2_disease_improvement.pdf}
\caption{Disease recognition improvement distribution. Point size represents training data proportion. Green points show >15\% gains, red points <-15\% declines.}
\label{fig:disease_improvement}
\end{figure}
%%%%%%%%%%%%%%%%%%%%%%%%%%%%%%%%%%%%%%%%%%%%%%%%%%%%%%%

\subsection{Analysis of Fine-Grained Disease Recognition}
Figure~\ref{fig:disease_improvement} visualizes disease recognition improvements. Diseases achieving >20\% gains cluster in the 3-8\% frequency range with moderate SFT baselines (45-65\%)—a "sweet spot" where sufficient gradient signal enables stable advantage estimation (Equation~\ref{eq:advantage}) and GRPO's exploration addresses visual complexity. Below 3\%, high-variance estimates cause erratic updates; above 8\%, SFT saturation limits further gains. Catastrophic forgetting in rare diseases. Four low-frequency diseases (<2\%) experience >20\% declines as high-frequency classes (>5\%) dominate batch sampling, overwriting their representations through gradient competition. This reflects a fundamental limitation of frequency-agnostic reward optimization—the KL term (Equation~\ref{eq:grpo_objective}) preserves overall stability but cannot prevent intra-class competition when gradient magnitudes differ by orders of magnitude.


%%%%%%%%%%%%%%%%%%%%%%%%%%%%%%%%%%%%%%%%%%%%%%%%%%%%%%%%
\begin{figure}[!t]
\centering
\includegraphics[width=\columnwidth]{figure4_complexity_comparison.pdf}
\caption{Ablation study on Disease Knowledge QA (20 samples). GRPO (no explicit reasoning, purple) provides baseline gains (+4\% to +28\%), while Reasoning-Enhanced GRPO (red) achieves 2.2× amplification on complex reasoning (+61\%).}
\label{fig:complexity}
\end{figure}
%%%%%%%%%%%%%%%%%%%%%%%%%%%%%%%%%%%%%%%%%%%%%%%%%%%%%%%%

%%%%%%%%%%%%%%%%%%%%%%%%%%%%%%%%%%%%%%%%%%%%%%%%%%%%%%%%
\begin{figure}[!t]
\centering
\includegraphics[width=\columnwidth]{figure5_cot-image.png}
\caption{
%Explicit reasoning example. Reasoning-Enhanced GRPO (top) generates structured \texttt{<think> ... </think>} reasoning with specific parameters (8.0/10) vs GRPO's answer lacking operational details (7.0/10).
A comparison of diagnostic reasoning. Our Reasoning-Enhanced GRPO (top) produces structured explanations with actionable details, while standard GRPO (bottom) provides minimal operational guidance.}
\label{fig:cot_example}
\end{figure}
%%%%%%%%%%%%%%%%%%%%%%%%%%%%%%%%%%%%%%%%%%%%%%%%%%%%%%%% 
\subsection{Analysis of Reasoning Capability}
\paragraph{ Scaling of Reasoning Benefits with Task Complexity}

The contribution of explicit reasoning scales dramatically with the complexity of the task. As shown in Figure~\ref{fig:complexity}, our analysis on the Disease Knowledge QA subset reveals a clear pattern: for highly complex, multi-domain questions, the improvement offered by standard GRPO remains at a robust +28\%. In contrast, our Reasoning-Enhanced GRPO framework achieves a remarkable +61\% gain, representing a 2.2× amplification of the performance improvement. This amplification results from the complementary roles of the two components: GRPO explores robust response patterns, while explicit reasoning chains provide the necessary scaffolding for multi-step problem-solving. This confirms that explicit reasoning fundamentally enhances capability, not just interpretability, for complex diagnostics.

\paragraph{Case Study: Qualitative Analysis of Explicit Reasoning Output}

A qualitative case study, illustrated in Figure~\ref{fig:cot_example}, highlights the practical advantage of generating structured reasoning. The output from our Reasoning-Enhanced GRPO model achieves a higher quality score (8.0/10) by producing a detailed, actionable reasoning chain.
%This includes specific operational parameters such as recommended solution temperature (56°C), treatment duration (5 minutes), chemical concentration (1\% KMnO₄), and precise dilution ratios (800-1200×).
In comparison, the response from the standard GRPO model, while still providing useful guidance (scoring 7.0/10), lacks these precise quantitative details. This omission limits its direct applicability for end-users, such as farmers or agricultural technicians, who require exact instructions for field implementation. The case demonstrates that our framework's reasoning synthesis not only boosts performance metrics but also generates outputs with significantly higher practical utility and operational specificity.




\section{Conclusion}
In this work, we introduce \textbf{Agri-R1, the first framework that integrates automated reasoning synthesis with GRPO for agricultural VQA task.} This approach pioneers a shift from opaque predictions to transparent, step-by-step diagnostic reasoning. Our contributions are twofold: (1) Scalable Reasoning Data Generation: We demonstrate how to automatically construct high-quality reasoning data without expert annotation, directly addressing the primary bottleneck for scaling interpretable AI in agriculture. (2) Domain-Specific Reward Design: We propose a novel domain-aware fuzzy-matching reward functions that effectively handle the linguistic diversity of open-ended agricultural responses—a challenge not fully addressed by RL systems from other domains. Empirical results confirm significant advantages: our framework enables superior cross-domain generalization, surpassing supervised methods that suffer from catastrophic performance drops. Crucially, our compact model achieves results competitive with larger baselines, proving that strategic learning, not merely parameter scale, is key for developing reasoning capabilities—an essential insight for deployment in resource-constrained environments.

We establish a new paradigm where explicit reasoning is both a performance driver and a practical necessity for building trust in real-world agricultural AI. Future work will focus on staged training protocols, frequency-aware optimization to mitigate class imbalance, and incorporating temporal modeling for dynamic disease progression.


% We introduce Agri-R1, the first framework integrating automated reasoning synthesis with Group Relative Policy Optimization for agricultural disease diagnosis, pioneering a paradigm shift from black-box predictions to interpretable, step-by-step diagnostic reasoning. This work delivers two foundational innovations that redefine agricultural AI: (1) automated construction of high-quality reasoning data without expert annotation, eliminating the primary scalability bottleneck; and (2) domain-aware fuzzy-matching rewards addressing linguistic diversity in open-ended agricultural responses—a critical challenge unaddressed by existing medical or mathematical RL systems.

% Empirical validation demonstrates transformative advantages across three dimensions. Reasoning-guided optimization enables robust cross-domain generalization, maintaining strong performance where supervised methods exhibit catastrophic collapse. Our compact model achieves competitive results with larger baselines, proving reasoning capability emerges from learning strategy rather than parameter scale—a paradigm-shifting insight for resource-constrained deployment. Explicit reasoning scales dramatically with task complexity, delivering substantial improvements on multi-step diagnostics integrating symptoms, context, and disease etiology.

% These findings establish a new agricultural AI paradigm, positioning interpretable reasoning as both a performance driver and practical necessity for farmer trust in real-world deployment. Future work will advance this framework through staged training, frequency-aware optimization, and temporal disease modeling.




\clearpage
\section*{Limitations}
Agri-R1 pioneers GRPO application to open-ended agricultural VQA, achieving robust generalization and interpretable reasoning without expert-annotated data. However, three technical limitations emerge. First, frequency-induced gradient competition causes rare-class representation degradation. Second, direct GRPO optimization without supervised warm-up results in format compliance issues. Third, the framework lacks temporal disease modeling. Future work will systematically address these through: (1) frequency-aware GRPO with curriculum learning prioritizing underrepresented classes; (2) staged training pipelines establishing structured output patterns via SFT before reward-guided optimization; and (3) recurrent vision-language architectures incorporating multi-temporal observations and environmental context. These limitations point to critical research directions but do not diminish the framework's paradigm-shifting contributions.


%\clearpage
\bibliography{custom,anthology-1,anthology-2}
% ,anthology-1,anthology-2

\appendix
% \section{Appendix}
% \label{sec:appendix}
%This appendix provides additional implementation details, including automated reasoning synthesis pipeline, reward function specifications, and training configurations.

\section{Automated Reasoning Synthesis Pipeline}
\label{sec:appendix_reasoning}

\subsubsection{DeepSeek-VL2 Generation Configuration}

\noindent We employ DeepSeek-VL2 with the following configuration: temperature=0.3 (low temperature ensures deterministic reasoning), max\_tokens=1024 (sufficient for 3-4 reasoning steps), top\_p=0.9 (nucleus sampling), frequency\_penalty=0.0.

\noindent System Prompt Structure: (1) Core Requirements—Output in English with 3-4 explicit steps labeled "Step 1:", "Step 2:", etc.; distinguish question types (PREVENTION/CONTROL vs IDENTIFICATION); keep reasoning concise (80-120 words) and practical. (2) Question Type Guidelines—PREVENTION: disease analysis → cultural practices → chemical control → application methods; IDENTIFICATION: plant morphology → symptom description → disease pattern → diagnosis with confidence. (3) Output Format—JSON with two fields: \texttt{\{"think": "Step 1: ...", "answer": "Final diagnosis"\}}.

\subsubsection{GPT-4 Quality Evaluation Rubric}

\noindent Generated outputs are evaluated using GPT-4 (gpt-4-0125-preview) with a 10-point rubric:

\vspace{3pt}
\begin{table}[t]
\centering
\small
\begin{tabular}{l@{\hspace{3pt}}c@{\hspace{3pt}}p{4.5cm}}
\hline
\textbf{Criterion} & \textbf{Score} & \textbf{Evaluation Focus} \\
\hline
Accuracy & 0-2 & Correct plant, disease identification; no hallucination \\[3pt]
Completeness & 0-2 & Key elements: plant, symptoms, disease \\[3pt]
Detail & 0-2 & Measurements, colors, distribution \\[3pt]
Relevance & 0-2 & Diagnosis-relevant; no redundancy \\[3pt]
Clarity & 0-2 & Professional terms; logical flow \\[3pt]
\hline
\textbf{Total} & \textbf{0-10} & \textbf{Threshold: $\tau$ = 8.0/10.0} \\
\hline
\end{tabular}
\caption{GPT-4 quality evaluation rubric for explicit reasoning outputs.}
\label{tab:cot_quality}
\end{table}
\vspace{2pt}

\noindent Filtering Process: Outputs scoring below $\tau$ = 8.0 trigger feedback-guided regeneration, where GPT-4's critique is appended to the original prompt. Maximum 2 regeneration attempts per sample. Samples failing after 2 attempts are excluded from the training set.

\subsubsection{Representative Reasoning Examples}

\noindent Example 1 (Disease Identification | Score: 8.5/10.0):

\noindent\textit{Question}: What is the content of this picture?

\noindent\textit{Ground Truth}: This image shows an apple leaf affected by Alternaria Blotch.

\noindent\textit{Generated Reasoning}:

\vspace{2pt}
\begin{footnotesize}
\begin{verbatim}
<think>Step 1: Identify plant - leaf ovate with
serrated margins and pinnate venation, consistent
with apple (Malus domestica). Step 2: Describe
symptoms - multiple circular brown lesions (2-5mm)
with yellowish halos scattered across leaf surface.
Step 3: Assess distribution - lesions cover ~20% of
visible area, no coalescence. Step 4: Preliminary
diagnosis - Alternaria Blotch caused by Alternaria
alternata; medium confidence based on lesion
morphology.</think>
<answer>Apple Alternaria Blotch</answer>
\end{verbatim}
\end{footnotesize}
\vspace{2pt}

\noindent Example 2 (Prevention/Control | Score: 9.0/10.0):

\noindent\textit{Question}: How to prevent and control tomato early blight?

\noindent\textit{Ground Truth}: Use crop rotation, remove infected debris, apply fungicides (chlorothalonil, mancozeb) every 7-10 days.

\noindent\textit{Generated Reasoning}:

\vspace{2pt}
\begin{footnotesize}
\begin{verbatim}
<think>Step 1: Disease analysis - Early blight
(Alternaria solani) overwinters in plant debris and
spreads via water splash. Step 2: Cultural practices
- 2-year rotation with non-Solanaceae crops; remove
infected leaves; mulch to prevent soil splash;
improve spacing for air circulation. Step 3:
Chemical control - Apply protectant fungicides
(chlorothalonil 1.5g/L, mancozeb 2g/L) at first
symptom; repeat every 7-10 days. Step 4:
Application - Spray thoroughly covering lower
leaves; use protective gear; avoid overhead
irrigation.</think>
<answer>Crop rotation, sanitation,
fungicide spray every 7-10 days</answer>
\end{verbatim}
\end{footnotesize}
\vspace{2pt}

\section{Reward Function Design}
\label{sec:appendix_reward}

\subsubsection{Domain Vocabulary Construction}
\label{sec:appendix_vocabulary}

To handle linguistic diversity in open-ended responses, we construct domain-specific vocabularies with comprehensive synonym recognition. Each plant/disease includes common names, scientific names, and colloquial variations extracted from CDDMBench annotations and agricultural databases.

\paragraph{Plant Variations (Complete List).}

Table~\ref{tab:plant_vocab_full} presents the complete plant vocabulary covering 15 major agricultural crops. Each entry includes the canonical name along with 5-8 recognized variations, encompassing scientific nomenclature (e.g., \emph{Solanum lycopersicum}), common names (e.g., ``tomato plant''), and colloquial terms (e.g., ``nightshade''). This comprehensive coverage enables robust plant identification despite lexical variation in model-generated responses.

\begin{table*}[t]
\centering
\small
\setlength{\tabcolsep}{4pt}
\begin{tabular}{lp{13.5cm}}
\hline
\textbf{Plant} & \textbf{Recognized Variations} \\
\hline
Tomato & tomato, tomato plant, tomatoes, solanum lycopersicum, lycopersicon esculentum, nightshade, tomato leaf, tomato crop \\
Potato & potato, potato plant, potatoes, solanum tuberosum, white potato, irish potato, potato tuber, potato crop \\
Corn & corn, corn plant, maize, zea mays, sweet corn, field corn, corn leaf, maize plant \\
Apple & apple, apple tree, malus domestica, apple crop, apple leaf, apple plant \\
Grape & grape, grapevine, vitis vinifera, grape plant, vineyard, grape leaf \\
Wheat & wheat, wheat plant, triticum aestivum, wheat crop, wheat leaf \\
Rice & rice, rice plant, oryza sativa, rice crop, paddy rice \\
Soybean & soybean, soy plant, glycine max, soya bean, soy crop \\
Bell Pepper & bell pepper, pepper plant, capsicum annuum, sweet pepper, pepper crop \\
Cherry & cherry, cherry tree, prunus avium, sweet cherry, cherry plant \\
Peach & peach, peach tree, prunus persica, peach crop \\
Strawberry & strawberry, strawberry plant, fragaria, strawberry crop \\
Blueberry & blueberry, blueberry plant, vaccinium, blueberry crop \\
Raspberry & raspberry, raspberry plant, rubus, raspberry crop \\
Pumpkin & pumpkin, pumpkin plant, cucurbita, pumpkin crop \\
\hline
\end{tabular}
\caption{Complete plant vocabulary (15 crops) with scientific and common name variations for fuzzy keyword matching.}
\label{tab:plant_vocab_full}
\end{table*}

\paragraph{Disease Variations (Complete List).}

Table~\ref{tab:disease_vocab_full} presents the complete disease vocabulary covering 20 disease categories plus healthy status. Each disease entry includes 3-6 recognized variations, incorporating pathogen-based names (e.g., \emph{Alternaria solani} for early blight), symptom-based descriptors (e.g., ``target spot''), and disease-type variations (e.g., ``alternaria leaf spot''). The ``healthy'' category requires exact matching to avoid false positives. This multi-faceted synonym structure handles the inherent ambiguity in agricultural disease nomenclature.

\begin{table*}[t]
\centering
\small
\setlength{\tabcolsep}{4pt}
\begin{tabular}{lp{12cm}}
\hline
\textbf{Disease} & \textbf{Recognized Variations} \\
\hline
Early Blight & early blight, alternaria solani, alternaria, target spot, alternaria leaf spot, early leaf blight \\
Late Blight & late blight, phytophthora infestans, phytophthora, oomycete disease, late leaf blight \\
Powdery Mildew & powdery mildew, erysiphales, white powdery coating, mildew, powdery fungus \\
Septoria Leaf Spot & septoria leaf spot, septoria, leaf spot disease, septoria blight \\
Mosaic Virus & mosaic virus, viral mosaic, mosaic disease, virus infection, viral disease, mosaic pattern \\
Leaf Mold & leaf mold, fulvia fulva, tomato leaf mold, fungal leaf mold, leaf mould \\
Bacterial Spot & bacterial spot, bacterial disease, bacterial leaf spot, bacteria infection, bacterial blight \\
Yellow Leaf Curl Virus & yellow leaf curl virus, ylcv, leaf curl virus, yellow leaf curl, viral leaf curl, tomato yellow leaf curl \\
Spider Mites & spider mites, mite damage, mite infestation, two-spotted spider mite \\
Target Spot & target spot, corynespora cassiicola, concentric lesions, target leaf spot \\
Leaf Rust & leaf rust, rust disease, rust fungus, leaf rust disease \\
Common Rust & common rust, corn rust, puccinia sorghi, maize rust \\
Northern Leaf Blight & northern leaf blight, turcicum leaf blight, leaf blight, northern corn leaf blight \\
Gray Leaf Spot & gray leaf spot, grey leaf spot, cercospora, gray spot \\
Leaf Scorch & leaf scorch, marginal leaf burn, leaf tip burn, scorch \\
Healthy & healthy, no disease, disease-free, normal plant, no symptoms, healthy plant, uninfected \\
Black Rot & black rot, rot disease, rotting, fungal rot, black root rot \\
Apple Scab & apple scab, scab disease, venturia inaequalis, scab \\
Alternaria Blotch & alternaria blotch, alternaria, blotch disease, alternaria leaf blotch \\
Leaf Blight & leaf blight, blight disease, blight, leaf blight disease \\
\hline
\end{tabular}
\caption{Complete disease vocabulary (20 diseases) with scientific names, common names, and pathogen-based variations.}
\label{tab:disease_vocab_full}
\end{table*}

\paragraph{Treatment Keywords (4 Categories).}

\noindent (1) Pesticides [0.6]: fungicide, copper, chlorothalonil, mancozeb, metalaxyl, azoxystrobin, propiconazole, captan, thiophanate, benomyl, bordeaux mixture, wettable powder.

\noindent (2) Cultural Practices [0.5]: crop rotation, air circulation, spacing, debris removal, resistant varieties, drainage, mulching, pruning, sanitation.

\noindent (3) Application Methods [0.5]: spray, application, protective gear, dosage, dilution, foliar application.

\noindent (4) Application Timing [0.4]: timing, early stage, first sign, onset, every 7-14 days, repeat, preventive application.

\subsubsection{Fuzzy Matching Implementation}
\label{sec:fuzzy_matching}

Unlike binary rewards in closed-set medical VQA, our fuzzy matching handles lexical variation through a five-tier scoring system designed specifically for agricultural terminology:

\begin{table}[h!]
\centering
\small
\setlength{\tabcolsep}{3pt}
\begin{tabular}{clp{5cm}}
\hline
\textbf{Tier} & \textbf{Score} & \textbf{Matching Criteria} \\
\hline
1 & 1.0 & Exact match: synonym from vocabulary \\
2 & 0.85 & High quality: multi-word term missing 1 word (e.g., "early" for "early blight") \\
3 & 0.7 & Partial: keyword stem matching (first 6 characters) \\
4 & 0.5 & Keyword: core words present (word length >3) \\
5 & 0.25 & Weak relevance: related terms (blight $\leftrightarrow$ disease/infection) \\
 $0$ & $0.0$ & No match \\
\hline
\end{tabular}
\caption{Five-tier fuzzy matching scoring system for agricultural terminology.}
\label{tab:fuzzy_matching}
\end{table}

\noindent Calculation Examples:

\vspace{2pt}
\noindent\textit{Example 1 (High-quality match)}:

\vspace{1pt}
\begin{footnotesize}
\begin{verbatim}
Reference: "Tomato Early Blight (Alternaria solani)"
Generated: "tomato plant with alternaria leaf spot"
Plant: 0.8 × 1.0 = 0.80 (exact: "tomato")
Disease: 1.2 × 0.85 = 1.02
  (high-quality: "alternaria" matches
  "early blight", missing "early")
Total: 1.82/2.0
\end{verbatim}
\end{footnotesize}
\vspace{2pt}

\noindent\textit{Example 2 (Perfect match)}:

\vspace{1pt}
\begin{footnotesize}
\begin{verbatim}
Reference: "Apple Powdery Mildew"
Generated: "apple tree with white powdery coating"
Plant: 0.8 × 1.0 = 0.80 (exact: "apple")
Disease: 1.2 × 1.0 = 1.20 (exact: "white powdery
  coating" in synonym list)
Total: 2.0/2.0 (perfect match)
\end{verbatim}
\end{footnotesize}
\vspace{2pt}

\noindent\textit{Example 3 (Weak match)}:

\vspace{1pt}
\begin{footnotesize}
\begin{verbatim}
Reference: "Tomato Bacterial Spot"
Generated: "tomato with bacterial infection"
Plant: 0.8 × 1.0 = 0.80 (exact: "tomato")
Disease: 1.2 × 0.5 = 0.60 (keyword: "bacterial"
  matches, but missing "spot")
Total: 1.40/2.0
\end{verbatim}
\end{footnotesize}
\vspace{2pt}

\subsubsection{Reward Component Specifications}

Our three-component reward function (Equation~\ref{eq:total_reward}) assigns weights based on diagnostic importance: Format (17\%) ensures structured output, Answer Keyword (67\%) directly measures diagnostic accuracy, and Reasoning (17\%) encourages logical coherence.

\paragraph{(1) Format Reward [0, 0.5] — 17\%.}

\noindent The format reward evaluates structural compliance and output quality through five sub-components. Table~\ref{tab:format_reward} provides detailed scoring criteria based on CDDMBench dataset statistics (mean reasoning length: 487 characters, mean answer length: 69 characters).

\begin{table*}[t]
\centering
\small
\setlength{\tabcolsep}{6pt}
\begin{tabular}{lcp{11cm}}
\hline
\textbf{Sub-component} & \textbf{Score} & \textbf{Evaluation Criteria} \\
\hline
Basic Structure & 0.15 & Must have both \texttt{<think> ... </think>} and \texttt{<answer> ... </answer>} tags; penalize if either missing \\
Step Structure & 0.15 & Number of explicit steps: $\geq$4 steps (0.15), 3 steps (0.12), 2 steps (0.08), 1 step (0.03) \\
Step Content Quality & 0.10 & Each step $\geq$30 characters: 4 valid steps (0.10), 3 steps (0.08), 2 steps (0.05), <2 steps (0.0) \\
Think Length & 0.05 & Optimal range 150-800 chars (0.05); acceptable 100-1000 chars (0.03); minimal $\geq$80 chars (0.01) \\
Answer Quality & 0.05 & Optimal range 15-200 chars (0.05); acceptable 10-300 chars (0.03); minimal $\geq$5 chars (0.01) \\
\hline
\textbf{Total} & \textbf{0.50} & \textbf{Dataset statistics: Mean Think=487 chars (SD=156), Mean Answer=69 chars (SD=28)} \\
\hline
\end{tabular}
\caption{Format reward breakdown (5 sub-components) with detailed scoring criteria based on CDDMBench statistics.}
\label{tab:format_reward}
\end{table*}

\paragraph{(2) Answer Keyword Reward [0, 2.0] — 67\%.}

\noindent The answer keyword reward dynamically evaluates diagnostic accuracy based on question type. For diagnosis questions, we employ weighted dual matching with fuzzy scoring (plant weight=0.8, disease weight=1.2). For treatment questions, we match against four method categories with tiered keyword counting.

\noindent Dynamic Evaluation based on question type:

\noindent\textit{Diagnosis Questions}:
\begin{itemize}
\setlength{\itemsep}{0pt}
\item Plant Match [0.8]: Extract plant name from reference answer using regex patterns (\texttt{(tomato|potato|corn|...)}), apply fuzzy matching (Table~\ref{tab:fuzzy_matching}) against PLANT\_VARIATIONS (Table~\ref{tab:plant_vocab_full}). Score = $0.8 \times$ fuzzy\_score (max 0.8).
\item Disease Match [1.2]: Extract disease name from reference answer using disease vocabulary, apply fuzzy matching against DISEASE\_VARIATIONS (Table~\ref{tab:disease_vocab_full}). Special handling: "healthy" status requires exact match (fuzzy=1.0 only if exact). Score = $1.2 \times$ fuzzy\_score (max 1.2).
\end{itemize}

\noindent\textit{Treatment Questions}:
\begin{itemize}
\setlength{\itemsep}{0pt}
\item Pesticides [0.6]: Count keywords from 16-term pesticide list. Tiered scoring: $\geq$3 keywords (0.6), 2 keywords (0.45), 1 keyword (0.3), 0 keywords (0.0).
\item Cultural Practices [0.5]: Count keywords from 15-term cultural practice list. Tiered scoring: $\geq$3 (0.5), 2 (0.35), 1 (0.2), 0 (0.0).
\item Application Methods [0.5]: Count keywords from 11-term application method list. Tiered scoring: $\geq$3 (0.5), 2 (0.35), 1 (0.2), 0 (0.0).
\item Application Timing [0.4]: Count keywords from 13-term timing list. Tiered scoring: $\geq$3 (0.4), 2 (0.3), 1 (0.15), 0 (0.0).
\end{itemize}

\paragraph{(3) Reasoning Reward [0, 0.5] — 17\%.}

\noindent The reasoning reward evaluates the quality of Chain-of-Thought explanations through three dimensions. Table~\ref{tab:reasoning_reward} details the evaluation criteria for logical coherence, professional terminology usage, and diagnostic chain completeness.

\begin{table*}[t]
\centering
\small
\setlength{\tabcolsep}{6pt}
\begin{tabular}{lcp{11cm}}
\hline
\textbf{Dimension} & \textbf{Score} & \textbf{Evaluation Criteria} \\
\hline
Logical Coherence & 0.25 & Presence of causal patterns: "observe...because", "symptom...indicate", "characteristic...suggest"; Step connections via related keywords (e.g., Step 1 mentions "leaf" → Step 2 describes "lesion") \\
Professionalism & 0.15 & Use of context patterns: "pathogen...infect", "symptom...show", "diagnosis...based on", "lesion...circular/brown"; Agricultural terminology (chlorosis, necrosis, pustule) \\
Completeness & 0.10 & Full diagnostic chain present: Observation phase (0.40) + Analysis phase (0.35) + Conclusion phase (0.35); Keywords: observe/see/visible (observation), analyze/indicate/disease (analysis), conclude/control/treatment (conclusion) \\
\hline
\textbf{Total} & \textbf{0.50} & \textbf{Encourages structured, professional diagnostic reasoning with complete observation-analysis-conclusion flow} \\
\hline
\end{tabular}
\caption{Reasoning reward breakdown (3 dimensions) evaluating logical structure, domain terminology, and diagnostic completeness.}
\label{tab:reasoning_reward}
\end{table*}

\paragraph{Design Rationale.}

The 67\% weight on Answer Keyword directly measures diagnostic accuracy—the primary objective for agricultural VQA. Format (17\%) ensures structured, parsable output for interpretability and downstream applications. Reasoning (17\%) encourages logical coherence and professional terminology, improving model trustworthiness. This weighting fundamentally differs from binary-reward closed-set medical VQA systems~\cite{lai2025med,pan2025medvlm}, which cannot handle the linguistic diversity inherent in open-ended agricultural responses.

\subsubsection{GRPO Training System Prompt}

\noindent The following system prompt guides the model during GRPO training, as implemented in our training script \texttt{train\_grpo\_with\_cot.sh} (available in code repository). The prompt distinguishes between identification and prevention/control questions, providing structured guidelines for each task type:

\vspace{4pt}
\noindent\rule{\columnwidth}{0.4pt}
\vspace{3pt}

\begin{footnotesize}
\begin{verbatim}
You are a plant disease management expert.
Carefully analyze the given image and question,
following these guidelines:

## Core Requirements:
1. Output must be in English and structured
   into explicit steps labeled
   "Step 1: ... Step 2: ..."
2. For PREVENTION/CONTROL questions: Focus
   ONLY on method reasoning - DO NOT
   re-diagnose the disease
3. For IDENTIFICATION questions: Focus on
   visual evidence and diagnostic reasoning
4. Keep reasoning concise (80-120 words)
   and practical

## Question Type Guidelines:

### FOR PREVENTION/CONTROL METHODS QUESTIONS:
- Step 1: Analyze disease characteristics
  that influence control strategies
  (pathogen biology, transmission)
- Step 2: Recommend cultural/preventive
  practices based on disease biology
  (rotation, sanitation, spacing)
- Step 3: Outline chemical control timing
  and selection (fungicide types,
  application intervals)
- Step 4: Integrate application methods
  and safety precautions (spray coverage,
  protective gear)

### FOR DISEASE IDENTIFICATION QUESTIONS:
- Step 1: Plant identification based on
  morphological features (leaf shape,
  margins, venation pattern)
- Step 2: Symptom observation and
  description (lesion color, size,
  distribution, shape)
- Step 3: Disease pattern analysis
  (spatial distribution, temporal
  progression, environmental conditions)
- Step 4: Preliminary diagnosis with
  confidence level (pathogen
  identification, differential diagnosis)

## CRITICAL RULES:
- If question asks about CONTROL/MANAGEMENT/
  PREVENTION/TREATMENT/METHODS: Use
  PREVENTION/CONTROL guideline
- If question asks about IDENTIFICATION/
  WHAT/NAME/DISEASE: Use IDENTIFICATION
  guideline
- NEVER mix guidelines - choose one based
  on question type

## Output Format:
<think>Step 1: ... Step 2: ... Step 3: ...
Step 4: ...</think>
<answer>Your final answer here</answer>
\end{verbatim}
\end{footnotesize}

\vspace{3pt}
\noindent\rule{\columnwidth}{0.4pt}
\vspace{4pt}



\begin{table*}[t]
\centering
\footnotesize
\setlength{\tabcolsep}{3.5pt}
\renewcommand{\arraystretch}{1.1}
\begin{tabular}{llp{10cm}}
\hline
\textbf{Category} & \textbf{Parameter} & \textbf{Value / Description} \\
\hline
\multicolumn{3}{l}{\textit{\textbf{Hardware \& Infrastructure}}} \\
& GPUs & 4 $\times$ NVIDIA A800 80GB (NVLink interconnect, 600GB/s bandwidth) \\
& CUDA Version & 11.8.0 with cuDNN 8.7.0 \\
\hline
\multicolumn{3}{l}{\textit{\textbf{Model Configuration}}} \\
& Base Model & Qwen2.5-VL-3B-Instruct (vision-language multimodal model) \\
& Model Parameters & 3.09B total (vision encoder: 0.67B, LLM: 2.42B) \\
& Image Resolution & 384×384 pixels (resized from original) \\
& Attention Mechanism & Flash Attention 2 (memory-efficient, 2-4× speedup) \\
& Precision & BF16 (Brain Float 16) for training stability \\
\hline
\multicolumn{3}{l}{\textit{\textbf{Optimization Configuration}}} \\
& Optimizer & AdamW (weight decay=0.01, $\beta_1$=0.9, $\beta_2$=0.999) \\
& Learning Rate & 8$\times$10$^{-7}$ (empirically tuned for GRPO stability) \\
& LR Scheduler & Cosine annealing with 15\% warmup (562 steps warmup) \\
& Warmup Strategy & Linear warmup from 0 to peak LR over 562 steps \\
& Max Gradient Norm & 0.3 (gradient clipping prevents instability) \\
& Weight Decay & 0.01 (L2 regularization) \\
& Gradient Checkpointing & Enabled (reduces memory by 30-40\%) \\
\hline
\multicolumn{3}{l}{\textit{\textbf{Batch \& Parallelism}}} \\
& Train Micro Batch Size & 10 per GPU (max fitting in 80GB memory) \\
& Gradient Accumulation & 4 steps (effective batch size amplification) \\
& Effective Batch Size & 160 (10 samples/GPU $\times$ 4 GPUs $\times$ 4 accum steps) \\
\hline
\multicolumn{3}{l}{\textit{\textbf{DeepSpeed ZeRO-3 Configuration}}} \\
& ZeRO Stage & 3 (parameter partitioning across GPUs) \\
& Offload Optimizer & None (keep optimizer states on GPU for speed) \\
& Offload Parameters & None (all parameters remain on GPU) \\
& Overlap Communication & True (overlap gradient communication with computation) \\
& Contiguous Gradients & True (memory layout optimization) \\
& Sub-group Size & 5e8 (500M parameters per partition) \\
& Reduce Bucket Size & 2e8 (200M, gradient reduction bucket) \\
& Stage3 Prefetch Bucket Size & 2e8 (200M, parameter prefetch bucket) \\
& Stage3 Param Persistence Threshold & 1e5 (100K, params kept in GPU) \\
& Stage3 Max Live Parameters & 1e9 (1B, max params in GPU memory) \\
& Stage3 Max Reuse Distance & 1e9 (1B, parameter reuse distance) \\
& Stage3 Gather 16bit Weights & True (gather BF16 weights on model save) \\
& Round Robin Gradients & True (distribute gradients evenly) \\
\hline
\multicolumn{3}{l}{\textit{\textbf{GRPO Strategy}}} \\
& Num Generations (K) & 3 candidates per sample (balance exploration vs cost) \\
& Sampling Temperature & 0.7 (moderate diversity for candidate generation) \\
& Top-p Sampling & 0.9 (nucleus sampling during generation) \\
& KL Divergence Coefficient & Auto-tuned (observed range: 0.036--0.040) \\
& KL Target & 0.04 (controls policy deviation from reference) \\
& Advantage Normalization & Group-relative (per-sample normalization) \\
& Clip Range & 0.2 (PPO-style clipping for stability) \\
\hline
\multicolumn{3}{l}{\textit{\textbf{Training Schedule}}} \\
& Num Train Epochs & 3 epochs (planned), 2.42 epochs (actual completion) \\
& Max Steps & 3,750 steps (planned), 3,027 steps (actual) \\
& Save Strategy & Every 300 steps (15 checkpoints total) \\
& Logging Steps & Every 2 steps (fine-grained monitoring) \\
& Evaluation Strategy & No evaluation during training (offline evaluation on test set) \\
\hline
\multicolumn{3}{l}{\textit{\textbf{Performance Metrics}}} \\
& Training Time & 98 hours ($\approx$4.1 days) for 3,027 steps \\
& Memory Usage & 68GB per GPU (near capacity) \\
& Optimal Checkpoint & Step 1,800 (Epoch 1.44, best test accuracy) \\
& Final Test Performance & Crop: 92.58\%, Disease: 72.50\% (checkpoint-1800) \\
\hline
\multicolumn{3}{l}{\textit{\textbf{Reproducibility}}} \\
& PyTorch Version & 2.1.0 (CUDA 11.8 build) \\
& Transformers Version & 4.36.0 (with Qwen2.5-VL support) \\
& DeepSpeed Version & 0.12.3 \\
\hline
\end{tabular}
\caption{Core training configuration parameters. Includes hardware setup (4×A800 GPUs), model architecture (Qwen2.5-VL-3B-Instruct), optimization settings (AdamW, lr=8e-7, gradient clipping), batch parallelism (10×4×4=160), DeepSpeed ZeRO-3 memory balancing, GRPO strategy (K=3, KL=0.036-0.040), training schedule, and performance metrics. Non-essential environment parameters omitted for brevity.}
\label{tab:training_config}
\end{table*}

\section{Empirical Configuration Choices}
\label{sec:empirical_config}

Each parameter was tuned to address key challenges: memory constraints through ZeRO-3 partitioning, gradient instability through aggressive clipping, and convergence efficiency through conservative learning rates.

\paragraph{Batch Size and Memory Management.}
Our effective batch size of 160 uses \texttt{train\_micro\_batch\_size\_per\_gpu}=10 with \texttt{gradient\_accumulation\_steps}=4 across 4 GPUs. This prioritizes stability: 20 samples/GPU caused frequent OOM. The 10×4 strategy reduces memory peaks while maintaining gradient quality, achieving 78-80GB utilization per GPU.

\paragraph{GRPO Candidates and Token Length.}
We use K=3 candidates per sample, reduced from Med-R1's K=4. Agricultural VQA elicits longer responses (mean 487/69 characters) than medical MCQs. K=4 exceeded 80GB; K=3 reduces memory by 25\% while maintaining stable KL.

\paragraph{Gradient Clipping and Stability.}
The 0.3 gradient norm threshold (vs typical 1.0) prevents explosion from reward variance in open-ended responses. Without clipping, norms exceeded 10.0 at steps 200-300, causing NaN losses. The threshold clips $\sim$15\% of gradients, validated by smooth KL curves.

\paragraph{DeepSpeed ZeRO-3 for Memory Balancing.}
We employ ZeRO Stage 3 for balanced GPU utilization. ZeRO-2 shows severe imbalance: GPU 0 at 90\% while GPUs 1-3 at 20-30\%, as it partitions optimizer states but replicates parameters. ZeRO-3 partitions parameters across devices, achieving 80-85\% on all GPUs and enabling 25\% higher batch size (8$\rightarrow$10/GPU). Key settings: \texttt{sub\_group\_size}=5e8, \texttt{overlap\_comm}=true ($\sim$15\% speedup), disabled CPU offloading (sufficient GPU memory).

\paragraph{Learning Rate and Convergence.}
The conservative 8×10$^{-7}$ learning rate reflects GRPO's sensitivity to policy deviation. Higher rates (5×10$^{-6}$) caused KL to exceed 0.04 within 500 steps. Our rate maintains stable KL throughout 3,027 steps. The 15\% warmup (562 steps) prevents early instabilities.

\paragraph{Precision and Attention Optimization.}
BF16 provides better stability than FP16, preventing activation overflow. Flash Attention 2 reduces memory by 30-40\%, enabling the 10×4×4 batch configuration and achieving 95\% GPU utilization across the 98-hour training.

\end{document}